\subsection{The annular square and its integrated estimate}

\begin{proposition}[Square of the covector-twisted Dirac operator]
\label{prop:annular-spin-square}
\uses{prop:supplied-normal-spin-defect}
\lean{FirstOrderTwistor.metricNormalChartTwistedDiracSquare, FirstOrderTwistor.metricFrameTwistedCurvaturePotential_eq_scalar_add, FirstOrderTwistor.metricFrameTwistedCurvaturePotential_pairing_le_intrinsic, FirstOrderTwistor.annularCovariantHessian_commutator, FirstOrderTwistor.annularDiracSquare_eq_traceHessian_add_curvature}
\leanok
For a $C^2$ positive tensor-spinor $v$ in the supplied aligned normal spin
chart of a smooth Riemannian four-metric, the square is
\[
 D^-D^+v=-\sum_i(\nabla^2 v)_{ii}+\mathcal Ev,\qquad
 \mathcal E=\tfrac14\mathrm{Scal}_g+
    \sum_{i<j}c^-(e_i)c^+(e_j)R^{T^*}_{ij}.
\]
The absolute curvature pairing is at most
$96|\mathrm{Rm}_g|\,|v|^2$ for the four-index tensor norm.
In arbitrary annular coordinates the same Hessian contraction uses $g^{ij}$.
The coordinate formula takes the Clifford anticommutator, the raised-symbol
connection equation, and the torsion-free tangent coefficients as inputs.
\end{proposition}
\begin{proof}
\leanok
Differentiate the tensor product connection in both covector slots.
The second covariant derivatives split into their symmetric Clifford
contraction and a commutator. The dual connection acts by negative
transpose. Contract the spin curvature using Bianchi and the Ricci
trace to obtain the scalar quarter. Estimate each of the six oriented
Clifford pairs in the curvature pairing and identify the full frame
norm with the intrinsic tensor norm.
\end{proof}

\begin{proposition}[Compact tests and relative regularization in the metric volume]
\label{prop:annular-spin-regularization}
\uses{prop:annular-spin-square,prop:spin-relative-bochner}
\lean{FirstOrderTwistor.annularDirac_green, FirstOrderTwistor.annularWeightedGradient_divergence_eq_traceHessian, FirstOrderTwistor.annularMetric_rough_le, FirstOrderTwistor.annularMetric_weitzenbock_le, FirstOrderTwistor.annularCliffordFluxCompatible_of_leviCivita, FirstOrderTwistor.spinMetricAnnulus_unweighted, FirstOrderTwistor.metricHarmonicFamily_weylMass_le_hessian, FirstOrderTwistor.spinCoordinateRegularizedFamilyProduct_at_normalCenter, FirstOrderTwistor.metricNormalChartHarmonicBochnerValues, FirstOrderTwistor.spinMetricAnnulus_regularized, FirstOrderTwistor.spinRelativeBochner_bound_of_jointZero, FirstOrderTwistor.spinJointZero_physicalCutoffCost}
\leanok
On a measured annulus, suppose the harmonic family, local spin germs,
curvature norms and compact test currents are realized in the physical
metric as in the linked declarations. For $S=\sum_a|\psi_a|^2$,
$Q=\sum_a|\nabla\psi_a|^2$ and every $\epsilon>0$, put
$r_\epsilon=Q/(S+\epsilon)$. Then
\[
 \int\chi^2|W^+|^2\frac S{S+\epsilon}\,dV_g
 \le 816\int\left(\chi^2|\mathrm{Ric}|^2\frac S{S+\epsilon}
 +\chi^2|\mathrm{Rm}|r_\epsilon+\chi^2r_\epsilon^2+
 |d\chi|^2r_\epsilon\right)dV_g .
\]
The right side uses the actual four-dimensional volume throughout.
For finite relative costs and $Q=0$ almost everywhere on
$\{S=0\}$ in the cutoff support, the limit bounds the weighted energy
on $\{S>0\}$. A comparison with the original energy measure charges
the remaining weighted energy on $\{S=0\}$ separately.
\end{proposition}
\begin{proof}
\leanok
Apply the metric divergence theorem to the compact current formed by
pairing $\nabla_i v+c^-(e_i)D^+v$ with $v$. Equivalently, in a measured
chart integrate the density-weighted Dirac and raised-gradient fluxes.
This yields $\int|\nabla v|^2\le\int|D^+v|^2+
96\int|\mathrm{Rm}|\,|v|^2$. Set $v=\chi\nabla\psi_a/\sqrt{S+\epsilon}$.
Sum the exact Ricci forcing and the differentiated cutoff estimates.
The two product terms have the same squared norm, and the
projected alternation costs at most twice the squared Hessian.
The coefficient from the regularized calculation is $48+8\cdot96$.
The final restriction to the positive locus takes place after the
compactly supported integration on the original patch.
\end{proof}

\begin{proposition}[Moving Clifford frame and gradient at a common zero]
\label{prop:spin-frame-joint-zero}
\uses{prop:annular-spin-square,prop:annular-spin-regularization}
\lean{FirstOrderTwistor.chiralSpinCovectorConnection_inner_skew, FirstOrderTwistor.spinNormalGaugeConnection_skew_of_orthonormal, FirstOrderTwistor.raisedFrameCliffordConnection_of_movingFrame, FirstOrderTwistor.lowerFrameCliffordConnection_of_movingFrame, FirstOrderTwistor.gradedCovectorClifford_skew, FirstOrderTwistor.gradedCovectorClifford_anticommute, FirstOrderTwistor.annularTwistedCurvaturePotential_gradedCovectorPositive_bound_iff, FirstOrderTwistor.annularInverseMetricConnection_of_movingFrame, FirstOrderTwistor.inverseMetricFromFrame_eq_gramInverse, FirstOrderTwistor.euclidean_spinJointZero_gradientSq_ae_on_open, FirstOrderTwistor.spinMetric_jointZero_gradientSq_ae}
\leanok
A supplied orthonormal frame with coordinate metricity has a skew
connection on both tensor-spin chiralities. Its raised and lowered
Clifford symbols satisfy the covariant first-order equations.
The frame also yields the inverse metric $UU^\top=G^{-1}$ and its
coordinate derivative. The positive curvature pairing agrees with its value in the full
graded Hilbert fibre. On an open spin chart the derivative density
$Q$ vanishes almost everywhere on the common zero set of a finite
$C^1$ family for every measure absolutely continuous with respect to
the chart volume. On a compact physical subset of a single chart this
conclusion holds for the actual metric volume.
\end{proposition}
\begin{proof}
\leanok
Differentiate $U^\top G U=I$ and contract the two chiral spin lifts
with the moving frame. The full Clifford action has two off-diagonal
blocks. Their product preserves the positive summand and its norm.
For the level-set assertion, take derivatives along coordinate lines,
apply Fubini, and intersect the finitely many full-measure sets.
The connection annihilates a field where its value is zero.
\end{proof}

\begin{proposition}[A harmonic zero at the critical exponent]
\label{prop:flat-spin-relative-zero}
\uses{prop:annular-spin-regularization}
\lean{FirstOrderTwistor.flatRelativeZeroSpinor_dirac, FirstOrderTwistor.flatRelativeZeroSpinor_mass_pos_ae, FirstOrderTwistor.flatRelativeZeroCost_not_integrable_ball_containing_zero, FirstOrderTwistor.flatRelativeZeroCost_annulus_integral, FirstOrderTwistor.flatComplementaryCost_integrable_ball, FirstOrderTwistor.flatComplementaryCost_annulus_logLower, FirstOrderTwistor.flatComplementaryCost_ball_tendsto_atTop}
\leanok
In the parallel positive spin frame of oriented flat $\mathbb R^4$,
$\psi(x)=x_0+x_1 i-x_2 j-x_3 k$ is harmonic and has a single zero.
Its relative fourth-power density integrates to
$32\pi^2\log(b/a)$ on $a<|x|\le b$ for $0<a\le b$.
It is nonintegrable on every ball containing the origin. Adding a
constant nonzero harmonic spinor $h$ gives a finite cost on each
bounded ball. At each fixed radius $R>0$ its integral over $B_R(0)$
tends to infinity as $h\to0$ through nonzero values.
\end{proposition}
\begin{proof}
\leanok
The Dirac contraction of the four constant coefficients vanishes.
The squared mass is $|x|^2$ and the squared gradient is four.
The radial density is $16/|x|^4$ and the unit three-sphere has
area $2\pi^2$. With the second field present, the density becomes
$16/(|x|^2+|h|^2)^2$. On $|h|<|x|\le b$, it is at least $4/|x|^4$.
Radial integration gives the claims.
\end{proof}

\begin{proposition}[Square-integrable incoming inverse in the spectral model]
\label{prop:weak-spectral-incoming}
\uses{prop:spin-mode-transmission}
\lean{FirstOrderTwistor.forcedCylinderMode_memH1, FirstOrderTwistor.forcedCylinderMode_weak_energy_balance, FirstOrderTwistor.cylinderGraphIncomingEquiv, FirstOrderTwistor.cylinderSmoothSingleModes_dense, FirstOrderTwistor.cylinderSmoothGraph_norm_sq_comparison, FirstOrderTwistor.cylinderMode_twoTrace_le, FirstOrderTwistor.cylinderWeakSynthesis, FirstOrderTwistor.CylinderWeakModalSolution.continuousTrace_realization, FirstOrderTwistor.CylinderWeakModalSolution.realize, FirstOrderTwistor.CylinderWeakModalSolution.weak_derivative_of_realization, FirstOrderTwistor.cylinderGraphIncomingPerturbation}
\leanok
In a fixed signed spectral model on $[0,L]$, $L>0$, assume a
countable Hilbert eigenbasis and the unit-round gap $|\lambda|\ge3/2$.
Any square-summable family of $L^2$ forcing coefficients and incoming
half-order traces has a unique modewise absolutely continuous solution.
Its squared bulk first-order graph norm is at most $13/9$ times the
sum of the forcing and incoming squared norms. The signed incoming
map is an isometric isomorphism in the graph norm consisting of bulk
derivative energy and the weighted outgoing trace. For $L\ge1$
this norm and the bulk norm are uniformly equivalent.
With fixed incoming projectors and a trace-preserving perturbation
of norm $\delta<1$, the inverse norm is at most $1/(1-\delta)$
in the isometric graph norm.
\end{proposition}
\begin{proof}
\leanok
Integrate the forcing from the left for positive modes and the right
for negative modes. The integrating factor gives an absolutely
continuous representative and its weak derivative. Integrating the reconstructed Hilbert-valued velocity from the left trace produces a single continuous representative. Every modal coefficient agrees on the closed interval, so both endpoint values agree there. The signed energy
identity has the outgoing trace on the bulk side. Sum it and use
the gap. Finite differentiable mode sums are dense in the completed
graph, and continuous forcing is dense in the modal $L^2$ spaces.
Use Parseval to synthesize the time and tangential derivatives in the
supplied basis. A linear cutoff gives a two-ended trace inequality
with explicit inverse-length term. Apply the Neumann bound to the
completed incoming map.
\end{proof}
